\section{Analysis}
\label{sec:analysis}

This section examines where the system created value, where it failed,
and what the failures taught us about the structure of the fantasy
scoring problem. Throughout, MD$k$ denotes matchday $k$ of the group
stage and R32 the round of 32.

\subsection{Calibration and sources of value}
\label{sec:analysis-value}

Figure~\ref{fig:calibration} plots pre-match predicted points against
realized points for every scored player-round. Two properties of this
plot frame the rest of the section. First, the predictions carry
genuine rank information: the per-position backend routing selected on
held-out EPL data (Table~\ref{tab:holdout-rmse}) generalized to the
tournament, with the heuristic the safest default for defenders, the
Poisson model for goalkeepers, and the GBM for attacking positions, and
no backend materially contradicted its own player ranking from round to
round. Second, the scale is compressed: realized points span a far
wider range than predicted points. This compression is not a defect of
any single backend but a structural property of mean prediction under
low-count scoring events, and it motivates the variance analysis in
Section~\ref{sec:analysis-variance}.

Beyond point prediction, two components delivered value that is easy to
under-credit. The rolling Elo signal (Section~\ref{sec:data-elo})
corrected real fixture mispricings: when the static FIFA ranking placed
Norway behind Senegal but the post-MD1 rolling Elo ranked Norway ahead,
the heuristic and Poisson backends adjusted Norway players upward by
approximately one expected point. This is the canonical case for
deriving live signals from realized results rather than hand-maintaining
static snapshots.

The MILP optimizer contributed in a way that RMSE cannot measure: it
kept every submitted squad legal. Budget, country-cap, and transfer
quota constraints are non-trivial to satisfy by hand under pre-deadline
time pressure. The solver never produced an illegal squad, and it
surfaced budget infeasibilities immediately (for example, that the
double swap Lautaro Mart\'inez to Messi plus Dou\'e to Bellingham costs
\$2.0M more than the available budget) rather than after a manual
mistake.

\subsection{Failure modes}
\label{sec:analysis-failures}

Three failure modes recurred. The first is GBM blindness to non-EPL
stars. The GBM has zero training rows for Messi and scores him from
price alone, adjusted downward by rotation-conservative status
defaults: at the R32 deadline the heuristic predicted 8.5 points for
Messi while the GBM predicted 2.2. A manager following the GBM alone
would have missed the tournament's top scorer. The heuristic backend,
whose price and matchup terms do capture such players, compensated
within the ensemble, but the premium tier remained the GBM's weakest
segment.

The second is coarse rotation-risk modeling. Our rotation multiplier
was a hand-set per-country scalar (for example 0.55 for Germany at MD3,
after Germany had clinched). A single scalar compresses three distinct
questions: who starts, for how many minutes, and whether the
replacement matters for fantasy purposes. A per-player signal informed
by manager press conferences would have been materially better; no such
feed existed. Section~\ref{sec:analysis-minutes} describes the partial
remedy we shipped.

The third is cross-domain distribution shift. Any feature whose range
at training time differs from its range at inference time is a risk,
and for attacking positions the GBM consistently under-predicted
premium international players relative to their realized totals.
Feature candidates rejected on these grounds are reported in
Section~\ref{sec:negative-results}.

A related operational failure concerns transfer hits. Each transfer
beyond the free quota costs 3 points, roughly 5\% of a typical
group-stage round for our squads, and hits compound. A hit is justified
only when the incoming player is expected to outscore the outgoing one,
net of captaincy implications, by more than 3 points over the planning
horizon. The model's transfer-with-hit logic was systematically too
aggressive because it underestimated rotation: the MD3 hits each cost 3
points and returned roughly 1 in realized output.

\subsection{The goalkeeper paradox}
\label{sec:analysis-gk}

The tournament produced one finding we did not anticipate: the best
real-world goalkeepers were not the best fantasy goalkeepers, and the
model structurally failed to predict this. Through the group stage,
Emiliano (Dibu) Mart\'inez of Argentina, widely regarded as the world's
best goalkeeper, totaled 14 fantasy points over three matches, while
Ra\'ul Rangel of Mexico, a competent domestic-league starter, totaled
21, half as much again. Both teams conceded zero goals across the three
matches.

The gap is entirely a save-bonus effect. The scoring rules award a
goalkeeper 2 points for a 60-minute appearance, 5 for a clean sheet,
1 per three saves, 5 per penalty save, and $-1$ per goal conceded after
the first. The save bonus is therefore increasing in shots faced.
Argentina's defense suppressed opponent shots so effectively that Dibu
made roughly one to two saves per match and rarely reached a save-bonus
threshold; Mexico conceded several attempts on target per match, and
Rangel converted them into save bonuses while keeping the same clean
sheets. The best real-world goalkeeper minimizes goals conceded per
shot faced; the best fantasy goalkeeper faces enough shots to
accumulate save bonuses while still shutting the opponent out. These
are different optimization problems.

Our model could not see the difference. The Poisson backend encoded the
save bonus as a flat constant of 1.0 expected points per start,
implicitly treating save opportunities as exogenous to the fixture. The
heuristic backend has the complementary defect: its matchup factor
inflates predictions with own-team strength for every position, but for
goalkeepers a strong own team reduces save opportunities while
clean-sheet probability for an elite defense is already near its
ceiling. Notably, the optimizer still selected Rangel over Dibu for
R32, at \$3.9M versus \$5.0M, but for the wrong reason: with equal
assumed save bonuses the choice was decided by price, and a different
flat constant could as easily have preferred Dibu. The correct decision
was reached through a model artifact, which is to say, fragilely.

The paradox implies a concrete goalkeeper profile: a team that keeps
clean sheets, a moderate level of opponent attacking threat (enough
shots to earn bonuses, not enough to break the shutout), an established
first-choice role, and a low price that frees budget for attacking
premiums. This explicitly contradicts the naive heuristic of buying the
best real-world goalkeeper. More generally, the scoring rules create
non-monotonic relationships between real-world quality and fantasy
value. We observed the same pattern elsewhere: centre-backs in deep
defensive blocks accumulate defensive counts, wide players in
possession-dominant teams see their credit diffused across long
build-up sequences, and substitute forwards on strong teams are worth
nothing unless they play, although price-based estimates implicitly
assume minutes.

\subsection{Empirically revising the save-bonus term}
\label{sec:analysis-gk-formula}

The fix replaces the flat constant with a term linear in opponent
expected goals: the predicted save bonus is $0.50$ times opponent xG
for the fixture. The multiplier was chosen empirically. We swept
candidate multipliers and measured goalkeeper RMSE on the held-out EPL
set (Figure~\ref{fig:gk-sweep}): the error curve has its minimum at
0.50, is shallow between 0.30 and 0.70, and deteriorates on both sides.
Table~\ref{tab:gk-ab} reports the resulting A/B comparison. Goalkeeper
RMSE improves from \GkRmseVOne{} (v1, flat) to \GkRmseVTwo{} (v2,
tuned) on the EPL holdout, with a larger improvement on realized World
Cup goalkeeper rounds; the other positions are unchanged because the
term applies only to goalkeepers. The revised term also has the correct
limiting behavior: a goalkeeper whose opponent creates nothing earns no
expected save bonus, where v1 awarded a full point.

The revision shipped under a two-gate rule: v2 replaced v1 only because
both the EPL held-out error and the realized World Cup error improved.
One objection deserves a direct answer: tuning a multiplier on held-out
data leaks information into the reported error. We tuned a single
scalar on one data set and validated the choice on an independent one,
and the selection among a handful of candidate values leaks little
relative to the size of the evaluation samples in
Table~\ref{tab:gk-ab}. Methodologically this follows
\citet{benter1994computer}: coefficients on plausible signals are
estimated against outcomes, not assumed. Our own first attempt is the
cautionary counterpart. We initially derived the multiplier
theoretically from shot and save rates; the derivation was internally
consistent and empirically wrong, and we report it with the other
negative results in Section~\ref{sec:negative-results}.

\subsection{Minutes and the availability discount}
\label{sec:analysis-minutes}

A benched player scores at most the appearance floor regardless of form
or matchup, so minutes information dominates any refinement of the
point model. The planned team-news pipeline returned zero predicted
lineup records for the World Cup, so the signal had to come from data
we already held. The World Cup player table exposes one realized total
per completed round; a player who took the pitch banks at least the
appearance points, so participation is proxied by a positive round
score. The feature \texttt{start\_rate\_lag} is the trailing fraction
of a player's last three completed rounds with a positive score; on the
EPL side the same feature uses the true starts flag. Both variants are
shifted by one match and are therefore leak-free. The signal is real:
recently ever-present players scored roughly twice as many points on
average as players who had rarely featured, and blanked far less often
(\POSTFINAL{number: participation-bucket means, blank rates, and the
correlation of start rate with realized points}).

The obvious move, feeding the proxy to the GBM as a feature, made the
model worse. Configuration D of the walk-forward study
(Table~\ref{tab:walkforward}), which adds the participation proxy
together with a team defensive-form feature examined in
Section~\ref{sec:negative-results}, regressed the pooled leak-free RMSE
from \WalkforwardPooledC{} to \WalkforwardPooledD{}. Over four rounds
the proxy is noisy, and it is collinear with lagged form (a benched
player already shows low form), so it adds variance the trees cannot
exploit.

The correct layer is the optimizer, not the model. The shipped
availability discount scales a player's effective predicted points by
$0.5 + 0.5 \cdot \texttt{start\_rate\_lag}$, a contraction toward a
floor of one half rather than zero, because the proxy is imperfect and
a nailed starter can blank twice in a row. Replayed on the heuristic
backend across MD1 to R32, the discount lifted realized squad points
(\POSTFINAL{number: realized squad points with and without the
availability discount, MD1--R32}), and the improvement was stable for
any floor between 0.3 and 0.7. The discount also surfaces an explicit
rotation-risk flag in the human-facing report. The lift is modest, but
for an entry at the bottom of its league every point counts.

\subsection{Ceiling-aware captaincy}
\label{sec:analysis-captaincy}

Captaincy doubles one player's score and is the single highest-leverage
decision in the game, and it is where model and manager disagreed most.
In MD1 both chose Lautaro Mart\'inez, who returned one raw point:
agreement, wrong outcome. In MD2 the model preferred Mbapp\'e; the user
overrode to Olise, who returned 9 raw points, the winning call. In MD3
the model preferred Bellingham; the user captained Messi in a favorable
fixture for 7 raw points. In R32 the model preferred Messi and the user
played Demb\'el\'e as a differential. Where they disagreed and the
outcome is known, the human call won, which we attribute to two inputs
the model does not have: an ownership-aware differential strategy
(maximum expected points is not maximum rank improvement for a trailing
entry) and football-watcher information such as visible form and
midweek fatigue.

Across the completed rounds, no backend selected the round-optimal
captain (Section~\ref{sec:live-results}). The root cause is the
objective: the lineup solver chose the captain as the argmax of mean
predicted points, and mean-argmax is the wrong rule for a trailing
team, which needs the highest ceiling rather than the highest
expectation. The GBM already emits a 90th-percentile (q90) head that
the captain choice ignored. A leak-free per-round backtest (train on
rounds before $k$, predict round $k$, form the eleven by mean) compared
captain rules on raw captain points captured across MD2, MD3, and R32:
the pure-ceiling q90 rule captured roughly twice the points of the
shipped mean rule, with an equal-weight blend in between and a
hindsight oracle far above all three
(\POSTFINAL{number: raw captain points captured by mean-argmax, blend,
q90-argmax, and oracle over MD2--R32}). Doubled by captaincy, the gap
amounts to several realized points per round, in a league being lost by
less.

The shipped selector scores captain candidates by
$(1 - c) \cdot \text{mean} + c \cdot \text{q90}$, where the chase
parameter $c$ is the entry's standings percentile: a leader ($c = 0$)
captains the safe mean, a last-placed entry ($c = 1$) captains the
ceiling. The optimizer now overrides the solver's naive pick with this
selector, defaulting $c = 0.9$ for our position near the bottom. For
point-estimate backends with no q90 head the rule degrades to the mean
and nothing breaks.

\subsection{Scoring variance is the strategy}
\label{sec:analysis-variance}

The user observed an oscillation in round scores: one round in the 40s,
the next in the 80s. The realized starter totals (captain double
excluded) were 32, 75, 39, and 73 across MD1 to R32, so the pattern is
real. It is not, however, a controllable cycle. The decisive evidence
is that MD1 and MD2 used the identical eleven starters and scored 32
and then 75: same players, no transfers, output more than doubled.
Fantasy points are driven by goals and assists, which are low-count
events, and an eleven-player squad naturally swings by tens of points
with selection held constant.

A mean point model is structurally unable to reproduce this. Summing
the heuristic's per-round predictions for the actual starters yields a
nearly flat series whose round-to-round standard deviation is a
fraction of the realized one (\POSTFINAL{number: realized versus
predicted round-total standard deviations}), consistent with the
amplitude compression visible in Figure~\ref{fig:calibration}. A
partially predictable component does exist: realized round totals
correlate positively with the average fixture-strength gap of the
starters and, more strongly, with the model's own per-round squad
prediction (\POSTFINAL{number: correlations of realized round totals
with fixture-strength gap and with the model's per-round prediction}).
The high-scoring rounds were the easier-fixture rounds, and the model
ranked rounds usefully even while understating their amplitude, so the
per-round squad prediction is a usable ex-ante ranker of which upcoming
round will be the high one.

The same lens resolves the mid-tournament observation that a
random-pick rival was outscoring the model-driven entry. A single round
proves little: our 42-point MD3 (net of a transfer hit) sits in the low
tail of the same process that produced a 102-point MD2. The honest
comparison is cumulative. Through R32, the ensemble replay accumulated
\TotalEnsemble{} points and the user's actual entry \TotalUserActualNet{}
net points, against a random-squad baseline mean of
\TotalRandomBaselineMean{} (Table~\ref{tab:season-summary}); the
model-driven process is far ahead of chance while remaining behind a
strong human who reads team news.

For a trailing entry in the knockout rounds, the consequence is the
opposite of smoothing the swing: variance is an asset, because a
low-variance strategy locks in the deficit. Concretely: captain by
ceiling, which the optimizer now does by default at our standings;
time aggression toward the rounds the model ranks as fixture-favorable;
prefer low-ownership differentials in those rounds, since a template
captain who hauls converts no rank gain; and do not spend transfers to
reduce round-to-round variance, because for a chaser the variance is
the only thing that can win it.

\begin{figure}[t]
  \centering
  \includegraphics[width=0.85\textwidth]{fig_calibration.pdf}
  \caption{Pre-match predicted points versus realized points.
    \figsource{fifa\_fantasy.report.data.load\_calibration $\rightarrow$
    data/processed}}
  \label{fig:calibration}
\end{figure}

\begin{figure}[t]
  \centering
  \includegraphics[width=0.85\textwidth]{fig_gk_sweep.pdf}
  \caption{Goalkeeper save-bonus multiplier sweep: empirical RMSE
    across candidate multipliers.
    \figsource{scripts/gk\_formula\_ab.py $\rightarrow$
    data/evaluation/gk\_formula\_ab.json}}
  \label{fig:gk-sweep}
\end{figure}

\begin{table}[t]
  \centering
  % GENERATED by scripts/paper_tables.py -- do not edit
\begin{tabular}{lrrrr}
\toprule
Dataset & $n$ & RMSE v1 & RMSE v2 & $\Delta$RMSE \\
\midrule
epl holdout & 180 & 2.503 & 2.491 & -0.013 \\
epl holdout & 910 & 3.400 & 3.400 & 0.000 \\
epl holdout & 1228 & 4.371 & 4.371 & 0.000 \\
epl holdout & 368 & 4.560 & 4.560 & 0.000 \\
wc realised & 168 & 3.790 & 3.640 & -0.150 \\
\bottomrule
\end{tabular}

  \caption{Goalkeeper save-bonus formula A/B: v1 (flat) versus v2
    (empirically tuned 0.50 multiplier).
    \figsource{scripts/gk\_formula\_ab.py}}
  \label{tab:gk-ab}
\end{table}
